% Part:sets-functions-relations
% Chapter: sets
% Section: schroder-bernstein

\documentclass[../../../include/open-logic-section]{subfiles}

\begin{document}

\olfileid[ar]{sfr}{siz}{sb}

\olsection{مفهوم الحجم ومبرهنة شرودر–برنشتاين}

\begin{explain}
إليك حدسًا أوليًا: إذا كانت $A$ لا تزيد حجمًا على $B$، وكانت $B$ لا
تزيد حجمًا على $A$، فإن $A$ و$B$ متساويتان عددًا. والحق أنه لو كان
هذا الحدس \emph{خاطئًا} لما استطعنا، إلا بعسر، أن نسوّغ القول إن لمفهوم
التساوي العددي الذي عرّفناه أي صلة بمقارنة «أحجام» المجموعات!{} ومن حسن
الحظ أن الحدس صحيح؛ وتبرهن مبرهنة شرودر–برنشتاين صحته.
\end{explain}

\begin{thm}[شرودر–برنشتاين]
  \ollabel{thm:schroder-bernstein}
  إذا كان $\cardle{A}{B}$ و$\cardle{B}{A}$،
  فإن $\cardeq{A}{B}$.
\end{thm}

\begin{explain}
وبعبارة أخرى، إذا وجدت !!a{injection} من $A$ إلى~$B$، ووجدت
!!a{injection} من $B$ إلى~$A$، وجدت !!a{bijection} من $A$
إلى~$B$.

غير أن برهان هذه النتيجة \emph{صعب} إلى حد بعيد. والواقع أن كانتور،
وإن كان قد صرّح بالنتيجة، فإن آخرين هم الذين برهنوها.\footnote{للمزيد
عن تاريخها، انظر مثلًا \citet[ص~165--6]{Potter2004}.}
\oliflabeldef{sfr:cardinals:card-sb:sec}{لن نكون في موضع يتيح لنا
\emph{إثبات} مبرهنة شرودر–برنشتاين إلا في
\olref[sfr][cardinals][card-sb]{sec}.}{}%
أما الآن، فيمكنك (ويجب عليك) أن تسلّم بصحتها.

ومن حسن الحظ أن مبرهنة شرودر–برنشتاين \emph{صحيحة}، وهي تؤيد تفسيرنا
للعلاقتين اللتين عرّفناهما، أي
$\cardeq{A}{B}$ و$\cardle{A}{B}$، باعتبارهما متصلتين بمفهوم «الحجم».
وفوق ذلك، فإن مبرهنة شرودر–برنشتاين \emph{مفيدة} جدًا؛ إذ قد يصعب تصور
!!a{bijection} بين مجموعتين متساويتين عددًا. وتتيح لنا مبرهنة
شرودر–برنشتاين تفكيك المقارنة إلى مسألتين، فلا يلزمنا إلا التفكير في
!!a{injection} من المجموعة الأولى إلى الثانية، وفي عكس الاتجاه.
\end{explain}
\end{document}
